\documentclass[10pt,a4paper]{amsart}
\usepackage{amsmath,amssymb,amsthm,graphicx}
% Retained minimal preamble for the Stacks Project chapter flat.tex (More on Flatness).
% Verbatim excerpt of the chapter preamble: Xy-pic setup (native preamble lines 14-18) and the native
% theorem declarations with the mathematical macros (native preamble lines 154-217). Copied mechanically;
% nothing added, removed or redefined. The wrapper class and the amsmath/amssymb/amsthm/graphicx base
% packages are supplied by the assembler, so the class/externaldocument part of the native preamble is omitted.
% For commutative diagrams we use Xy-pic
\usepackage[all]{xy}

% We use 2cell for 2-commutative diagrams.
\xyoption{2cell}
% Theorem environments.
%
\theoremstyle{plain}
\newtheorem{theorem}[subsection]{Theorem}
\newtheorem{proposition}[subsection]{Proposition}
\newtheorem{lemma}[subsection]{Lemma}

\theoremstyle{definition}
\newtheorem{definition}[subsection]{Definition}
\newtheorem{example}[subsection]{Example}
\newtheorem{exercise}[subsection]{Exercise}
\newtheorem{situation}[subsection]{Situation}

\theoremstyle{remark}
\newtheorem{remark}[subsection]{Remark}
\newtheorem{remarks}[subsection]{Remarks}

\numberwithin{equation}{subsection}

% Macros
%
\def\lim{\mathop{\mathrm{lim}}\nolimits}
\def\colim{\mathop{\mathrm{colim}}\nolimits}
\def\Spec{\mathop{\mathrm{Spec}}}
\def\Hom{\mathop{\mathrm{Hom}}\nolimits}
\def\Ext{\mathop{\mathrm{Ext}}\nolimits}
\def\SheafHom{\mathop{\mathcal{H}\!\mathit{om}}\nolimits}
\def\SheafExt{\mathop{\mathcal{E}\!\mathit{xt}}\nolimits}
\def\Sch{\mathit{Sch}}
\def\Mor{\mathop{\mathrm{Mor}}\nolimits}
\def\Ob{\mathop{\mathrm{Ob}}\nolimits}
\def\Sh{\mathop{\mathit{Sh}}\nolimits}
\def\NL{\mathop{N\!L}\nolimits}
\def\CH{\mathop{\mathrm{CH}}\nolimits}
\def\proetale{{pro\text{-}\acute{e}tale}}
\def\etale{{\acute{e}tale}}
\def\QCoh{\mathit{QCoh}}
\def\Ker{\mathop{\mathrm{Ker}}}
\def\Im{\mathop{\mathrm{Im}}}
\def\Coker{\mathop{\mathrm{Coker}}}
\def\Coim{\mathop{\mathrm{Coim}}}

% Boxtimes
%
\DeclareMathSymbol{\boxtimes}{\mathbin}{AMSa}{"02}

%
% Macros for moduli stacks/spaces
%
\def\QCohstack{\mathcal{QC}\!\mathit{oh}}
\def\Cohstack{\mathcal{C}\!\mathit{oh}}
\def\Spacesstack{\mathcal{S}\!\mathit{paces}}
\def\Quotfunctor{\mathrm{Quot}}
\def\Hilbfunctor{\mathrm{Hilb}}
\def\Curvesstack{\mathcal{C}\!\mathit{urves}}
\def\Polarizedstack{\mathcal{P}\!\mathit{olarized}}
\def\Complexesstack{\mathcal{C}\!\mathit{omplexes}}
% \Pic is the operator that assigns to X its picard group, usage \Pic(X)
% \Picardstack_{X/B} denotes the Picard stack of X over B
% \Picardfunctor_{X/B} denotes the Picard functor of X over B
\def\Pic{\mathop{\mathrm{Pic}}\nolimits}
\def\Picardstack{\mathcal{P}\!\mathit{ic}}
\def\Picardfunctor{\mathrm{Pic}}
\def\Deformationcategory{\mathcal{D}\!\mathit{ef}}
\usepackage{cleveref}
\crefname{originaldefinition}{original definition}{original definitions}
\Crefname{originaldefinition}{Original definition}{Original definitions}
\makeatletter
\newlabel{definition-one-step-devissage}{{05H4}{source}{}{sourcepointer-1}{}}
\newlabel{definition-one-step-devissage@cref}{{[originaldefinition][0][]05H4}{[1][0][]source}{}{}{}}
\makeatother
\crefname{originallemma}{original lemma}{original lemmas}
\Crefname{originallemma}{Original lemma}{Original lemmas}
\makeatletter
\newlabel{algebra-lemma-NAK}{{00DV}{source}{}{sourcepointer-2}{}}
\newlabel{algebra-lemma-NAK@cref}{{[originallemma][0][]00DV}{[1][0][]source}{}{}{}}
\makeatother
\title{A finite free cover of the pushforward of a one step devissage with a strictly smaller cokernel on the fibre}
\author{The Stacks Project Authors}
\date{stacks-project:flat.tex:local-snapshot-20261011; GNU FDL 1.2 or later}
\begin{document}
\maketitle
\noindent\textit{Editorial note.} All delivered mathematical text is the author's own, copied from the pinned Stacks Project chapter file flat.tex; no statement or proof step was rewritten, reordered or removed.\par\medskip
\noindent\textit{Editorial note.} The preamble is the chapter's own retained minimal preamble (sources/stacks-flat/retained-preamble.tex), a verbatim excerpt of the native Xy-pic setup and of the native theorem declarations and mathematical macros; no macro was invented or redefined.\par\medskip
\noindent\textit{Editorial note.} The same-chapter referenced definition is located at its \textbackslash{}label line inside the pinned chapter file; the cross-chapter Nakayama lemma is located in the project's own tag list staged in the snapshot (sources/stacks-chow/tags.txt), at the verbatim line carrying its canonical tag.\par\medskip
\noindent\textit{Allowed standard prerequisites (editorial).}\begin{itemize}
\item One step devissage of a finite type module over a point. Same chapter, Definition environment labelled definition-one-step-devissage, source line 780 (Stacks tag 05H4). Supplies the data (Z, Y, i, pi, G) and the affineness of S and Y that the counted proof starts from, plus finite generation of pi\_*G; nothing of the construction of alpha is taken from it.
\item Nakayama's lemma. Algebra, Lemma labelled algebra-lemma-NAK (canonical Stacks tag 00DV); invoked at proof line 1178. Used once, to promote the vanishing N' tensor kappa(p) = 0 to the vanishing of the localisation N'\_p; the choice of r, of the basis x\_1, ..., x\_r and the dimension comparison are inside the counted span.
\end{itemize}\medskip
\ifdefined\hypertarget\hypertarget{sourcepointer-1}{}\fi
\noindent\textit{Source reference.} Original definition 05H4: sources/stacks-flat/pinned-original-source.tex, source lines 780--780. Same-chapter pointer to the definition named in the statement at line 1140; it fixes the hypothesis that S and Y are affine and that pi is finite, all of which the counted proof uses explicitly.\par\medskip
\ifdefined\hypertarget\hypertarget{sourcepointer-2}{}\fi
\noindent\textit{Source reference.} Original lemma 00DV: sources/stacks-chow/tags.txt, source lines 512--512. Cross-chapter standard prerequisite invoked once at proof line 1178; it supplies only the promotion of N' tensor kappa(p) = 0 to N'\_p = 0, while the construction of alpha and the geometric-irreducibility dimension comparison are retained in the counted proof.\par\medskip
\begin{lemma}
\label{lemma-existence-alpha}
Let $S$, $X$, $\mathcal{F}$, $s$ be as in
Definition \ref{definition-one-step-devissage}.
Let $(Z, Y, i, \pi, \mathcal{G})$ be a one step d\'evissage
of $\mathcal{F}/X/S$ over $s$.
Let $\xi \in Y_s$ be the (unique) generic point.
Then there exists an integer $r > 0$ and an $\mathcal{O}_Y$-module map
$$
\alpha : \mathcal{O}_Y^{\oplus r} \longrightarrow \pi_*\mathcal{G}
$$
such that
$$
\alpha :
\kappa(\xi)^{\oplus r}
\longrightarrow
(\pi_*\mathcal{G})_\xi \otimes_{\mathcal{O}_{Y, \xi}} \kappa(\xi)
$$
is an isomorphism. Moreover, in this case we have
$$
\dim(\text{Supp}(\Coker(\alpha)_s)) < \dim(\text{Supp}(\mathcal{F}_s)).
$$
\end{lemma}

\begin{proof}
By assumption the schemes $S$ and $Y$ are affine.
Write $S = \Spec(A)$ and $Y = \Spec(B)$.
As $\pi$ is finite the $\mathcal{O}_Y$-module $\pi_*\mathcal{G}$
is a finite type quasi-coherent $\mathcal{O}_Y$-module.
Hence $\pi_*\mathcal{G} = \widetilde{N}$ for some finite $B$-module $N$.
Let $\mathfrak p \subset B$ be the prime ideal corresponding to $\xi$.
To obtain $\alpha$ set
$r = \dim_{\kappa(\mathfrak p)} N \otimes_B \kappa(\mathfrak p)$
and pick $x_1, \ldots, x_r \in N$ which form a basis of
$N \otimes_B \kappa(\mathfrak p)$. Take $\alpha : B^{\oplus r} \to N$
to be the map given by the formula $\alpha(b_1, \ldots, b_r) = \sum b_ix_i$.
It is clear that
$\alpha : \kappa(\mathfrak p)^{\oplus r} \to N \otimes_B \kappa(\mathfrak p)$
is an isomorphism as desired. Finally, suppose $\alpha$ is any map with this
property. Then $N' = \Coker(\alpha)$ is a finite $B$-module
such that $N' \otimes \kappa(\mathfrak p) = 0$. By Nakayama's lemma
(Algebra, Lemma \ref{algebra-lemma-NAK})
we see that $N'_{\mathfrak p} = 0$. Since the fibre $Y_s$ is
geometrically irreducible of dimension $n$ with generic point $\xi$
and since we have just seen that $\xi$ is not in the support of
$\Coker(\alpha)$ the last assertion of the lemma holds.
\end{proof}

\end{document}
